%==============================================================================
\section{Introduction}
%==============================================================================
% Structure follows the supervisor's 5-paragraph / 9-sentence template:
% hook, gap, study (+RQs), contribution, roadmap.

\rev{Large Language Models (LLMs)} built on the Transformer architecture \cite{vaswani2017attention} now power knowledge-intensive applications\rev{, such as} question answering, search, and code assistance, yet their parametric knowledge is frozen at training time and they often hallucinate, producing fluent text that is factually wrong \cite{ji2023survey}. Retrieval-Augmented Generation (RAG) \cite{lewis2020rag} addresses this by retrieving documents from an external corpus and conditioning generation on them, and it has become a default pattern for deploying LLMs over private or fast-changing knowledge \cite{gao2024retrieval}.

Grounding generation in retrieved text creates a new dependency: the answer is only as trustworthy as the corpus behind it. Standard RAG assumes a benign knowledge base, and existing evaluation frameworks\rev{, such as} RAGAS \cite{es2024ragas} and ARES \cite{saadfalcon2024ares}\rev{,} measure faithfulness to the retrieved context rather than the integrity of that context. In open-world deployments, however, an adversary can insert malicious documents through knowledge poisoning \cite{zou2024poisoned}; injecting only five malicious passages per question into a corpus of millions ($\approx$0.0002\%) can drive attack success to roughly 90\%. A system that answers faithfully from a poisoned document therefore scores well under current metrics while emitting compromised output. We call this the \emph{Security-Reliability Gap}: high semantic relevance is treated as a proxy for truth, and the few available poisoning defenses are mostly offline corpus-cleaning steps with no online trust layer inside the inference loop.

This paper presents an \emph{Evaluation Agent} that \rev{closes this Security-Reliability Gap}. It acts as defensive middleware that screens retrieved context before generation and outputs an interpretable \emph{Trust Index} that fuses Natural Language Inference (NLI) factual verification \cite{maynez2020faithfulness}\cite{honovich2022true}, a five-signal poison detector \rev{(Section~III)}, and a cross-document consistency estimate. We evaluate the agent on two public benchmarks (TruthfulQA and FEVER) across three LLMs and four attack strategies, and then apply it to a software-engineering use case, a secure-coding assistant that retrieves from guidance based on the Open Worldwide Application Security Project (OWASP) Top~10 and the Common Weakness Enumeration (CWE). We study three research questions:
\begin{itemize}
  \item \textbf{RQ1.} How effectively does the Evaluation Agent detect misinformation and knowledge poisoning in retrieved context?
  \item \textbf{RQ2.} How does NLI-based factual verification change the trustworthiness of RAG output?
  \item \textbf{RQ3.} How resilient is the approach across LLMs, datasets, and attack strategies, including a secure-coding setting in the software-development lifecycle (SDLC)?
\end{itemize}
\rev{The RQs operationalize this gap: RQ1 targets detection, RQ2 the trust layer that NLI adds inside the inference loop, and RQ3 its robustness across generators, domains, and attacks; Section~VI answers each explicitly. Code, prompts, the attack generator, and all experimental artifacts are publicly available} \cite{repoartifacts}.

We find that the agent detects overt poisoning well (100\% recall on instruction injection, 91\% accuracy and 100\% precision on TruthfulQA mixed attacks) but that in-place edits\rev{, such as} entity swaps and subtle weakening\rev{,} stay near-undetectable, which is a limit of surface-signal detection rather than of tuning. We also show that NLI-based trust scoring depends on the generating LLM (generation style matters more than model size), is stable across repeated runs, and is invariant to retrieval depth, while a weaker FEVER result bounds external validity. These findings extend prior RAG evaluation by adding an online check on context integrity, mapping where such checks succeed and where world-knowledge verification is still needed.

%The remainder of the paper is organized as follows. Section~II reviews background and related work. Section~III presents the Evaluation Agent, threat model, and Trust Index. Section~IV describes the experimental design, Section~V reports results, Section~VI discusses implications and limitations, and Section~VII concludes.
